\documentclass[11pt,a4paper]{article}
\usepackage[T1]{fontenc}
\usepackage{lmodern,amsmath,amssymb,amsthm,microtype}
\usepackage[margin=27mm]{geometry}
\usepackage[colorlinks=true,linkcolor=blue,urlcolor=blue]{hyperref}
\newtheorem{theorem}{Theorem}[section]
\newtheorem{lemma}[theorem]{Lemma}
\newcommand{\Q}{\mathbb Q}
\DeclareMathOperator{\Odd}{Odd}
\title{An obstruction to multiplying six-contact functions}
\author{}\date{10 November 2026}
\begin{document}\maketitle
\begin{abstract}
The six rational contacts used in the weak-$abc$ construction admit
a rational involution which changes the sign of their product modulo
squares. We give positive rational scales of unbounded height for which
the product of the two contact functions has even valuation at every
finite prime. Consequently the odd radical of that product cannot by
itself control the scale height on the unrestricted rational domain.
This obstruction is independent of any $abc$ conjecture. It does not
assert the existence of witnesses for the remaining source equations
or for their prescribed eighth-power scale.
\end{abstract}

\section{The contact function and its involution}

For $u\in\Q^\times$, define its odd radical by
\[
             \Odd(u)=\prod_{\substack{p\text{ prime}\\v_p(u)\text{ odd}}}p.
\]
Here negative valuations are included. Thus $\Odd(u)=1$ if and only
if $u$ is a rational square or the negative of a rational square.
For $u=a/b$ in lowest terms with $b>0$, put
$h(u)=\log\max(|a|,b)$.

Let
\[
 \Theta=\{0,-1/2,-1/4,-2/5,-1,-1/3\}.
\]
Put
\[
 F(t)=\prod_{\theta\in\Theta}(t-\theta)
     =\frac{t(t+1)(2t+1)(3t+1)(4t+1)(5t+2)}{120},
\]
and define
\[
 R(t)=\frac{F(t)}{(t-1)^8},\qquad
 \phi(t)=-\frac{2t+1}{3t+2}.
\]
These are the six contacts and the contact function of the
weak-$abc$ construction in~\cite{Conditional}. The rational map
$\phi$ is an involution of the projective line. It exchanges the
contact pairs
\[
 \{0,-1/2\},\qquad\{-1/4,-2/5\},
\]
and fixes $-1$ and $-1/3$.

\begin{lemma}\label{lem:identity}
For every rational $t$ for which the displayed expressions are
defined,
\begin{equation}\label{eq:identity}
 F(\phi(t))=-\frac{F(t)}{(3t+2)^6},\qquad
 R(t)R(\phi(t))
 =-\left(\frac{F(t)(3t+2)}{(t-1)^4(5t+3)^4}\right)^2.
\end{equation}
In particular the second product has even valuation at every finite
prime whenever it is nonzero.
\end{lemma}
\begin{proof}
Write $N=-2t-1$ and $D=3t+2$. Then
\[
 \begin{aligned}
 N+D&=t+1,&2N+D&=-t,&3N+D&=-3t-1,\\
 4N+D&=-5t-2,&5N+2D&=-4t-1.
 \end{aligned}
\]
Consequently
\[
 N(N+D)(2N+D)(3N+D)(4N+D)(5N+2D)=-120F(t).
\]
Dividing by $120D^6$ proves the first identity. Also
$\phi(t)-1=-(5t+3)/(3t+2)$, whence
\[
 R(\phi(t))=-\frac{F(t)(3t+2)^2}{(5t+3)^8}.
\]
Multiplying by $R(t)$ proves the second identity. Every finite
valuation of the negative of a nonzero square is even.
\end{proof}

\section{Unbounded positive scales}

\begin{theorem}\label{thm:obstruction}
For every integer $m\ge5$, put
\[
 \lambda_m=-\frac1m,\qquad
 a_m=\frac{m(m-2)}{2m-3}.
\]
Then $a_m>1$, both contact functions below are defined and nonzero,
and
\[
 a_m\lambda_m=\phi(\lambda_m),\qquad
 \Odd\bigl(R(\lambda_m)R(a_m\lambda_m)\bigr)=1.
\]
Moreover $h(a_m)\longrightarrow\infty$ as $m\longrightarrow\infty$.
Hence no constants $C,C_0$ give an inequality
\begin{equation}\label{eq:impossible}
 h(a)\le C\log\Odd\bigl(R(\lambda)R(a\lambda)\bigr)+C_0
\end{equation}
for all positive rational $a>1$ and all rational $\lambda$ for
which the two contact functions are defined and nonzero.
\end{theorem}
\begin{proof}
Direct substitution gives
\[
 \phi(-1/m)=-\frac{m-2}{2m-3}=a_m\lambda_m,
 \qquad a_m-1=\frac{(m-1)(m-3)}{2m-3}>0.
\]
For $m\ge5$,
\[
 -\frac14<\lambda_m<0,
 \qquad -\frac12<\phi(\lambda_m)<-\frac25.
\]
Neither open interval contains a contact or the pole $1$ of $R$.
Thus the product is nonzero and defined. Lemma~\ref{lem:identity}
makes its odd radical equal to one. Finally
$a_m>(m-2)/2$, so $h(a_m)\ge\log a_m\longrightarrow\infty$.
The right-hand side of~\eqref{eq:impossible} would be the fixed
constant $C_0$ on this family, a contradiction.
\end{proof}

Removing any fixed finite set of primes from the definition of
$\Odd$ leaves this obstruction unchanged. Multiplication or division
of the product by a nonzero rational square also leaves it unchanged.

\section{Scope and exact verification}

The theorem concerns an unrestricted contact-height inequality.
In the weak-$abc$ defining source, the scale is specifically
$a=\beta^8+T(n)^8$, where $T(n)$ is the elliptic trace, and the
contact expressions are accompanied by other equations and guards.
The family above is not asserted to satisfy that scale condition or
those equations. Thus it does not give a counterexample to the
thirteen-variable theorem, or exclude every possible construction
with one pole predicate. It shows that a direct product argument
needs additional arithmetic information about the scale or the
common contact primes. The obstruction holds whether or not $abc$
is assumed.

The standard-library program
\begin{center}\small
\nolinkurl{computation/pole_parameters/check_six_contact_product.py}
\end{center}
checks the homogenized degree-six identity, the involution matrix,
the contact permutation, and the parameterized scale identities by
exact polynomial arithmetic. It also evaluates the resulting
negative-square identity for every integer $5\le m\le200$ with
exact fractions. The polynomial identities and the proof above
establish the infinite family; the finite evaluations are additional
checks, not the basis for that assertion.

\begin{thebibliography}{9}
\bibitem{Conditional}
\emph{Thirteen Variables from a Weak $abc$ Inequality},
companion paper, 9 November 2026,
\nolinkurl{conditional13_20261109_v16.tex}.
\end{thebibliography}
\end{document}
